\documentclass[12pt]{article}
\usepackage[spanish]{babel}
\usepackage{natbib}
\usepackage{url}
\usepackage[utf8x]{inputenc}
\usepackage{amsmath}
\usepackage{graphicx}
\usepackage{parskip}
\usepackage{fancyhdr}
\usepackage{vmargin}
\usepackage[ddmmyy]{datetime}
\usepackage{anyfontsize}
\usepackage{helvet}
\renewcommand{\familydefault}{phv}
\usepackage{xcolor}

% \setmarginsrb{leftmargin}{topmargin}{rightmargin}{bottommargin}%
%          {headheight}{headsep}{footheight}{footskip}
\setmarginsrb{2.5 cm}{2.5 cm}{2.5 cm}{2.00 cm}
           {1 cm}{1.25 cm}{1 cm}{1.5 cm}

\usepackage{color}
\usepackage{hyperref}
\hypersetup{
    colorlinks=true, % set true if you want colored links
    linktoc=all,     % set to all if you want sections and subsections linked
    linkcolor=blue,  % choose some color if you want links to stand out
    urlcolor=blue,
}

\usepackage{subfig}
\usepackage[labelformat=empty]{caption}

\usepackage{fancyvrb}
\fvset{xleftmargin=\mathindent}

% code blocks
\usepackage{verbatimbox}
\newenvironment{fullgrayverb}
{\verbbox}
{\endverbbox\par\colorbox{gray!25}{\parbox{\textwidth}{\theverbbox}}\par}

% inline code blocks
\usepackage{tcolorbox}
\newcommand\mystrut{\rule[-1pt]{0pt}{.8em}}
\newtcbox{\code}{on line, boxrule=0pt, boxsep=0pt, top=2pt,
left=2pt, bottom=2pt, right=2pt, colback=gray!25, colframe=white,
fontupper={\ttfamily\mystrut}}

%%%%%%%%%%%%%%%%%%%%%%%%%%%%%%%%%%%%%%%%%%%%%%%%%%%%%%%%%%%%%%%%%%%%%%%%%%%%%%%%
% DATOS GENERALES

\title{Números Primos}          % Titulo del trabajo.
\author{Martin J. Klöckner}     % Nombre y apellido
\newcommand{\padron}{123456}    % Padrón
\newcommand{\tpnumber}{1}       % Número de trabajo práctico
\date{\today}                   % Fecha (automática)

\makeatletter
\let\thetitle\@title
\let\theauthor\@author
\let\thedate\@date
\makeatother

\pagestyle{fancy}
\fancyhf{}
\rhead{\theauthor}
\lhead{\thetitle}
\cfoot{\thepage}

%%%%%%%%%%%%%%%%%%%%%%%%%%%%%%%%%%%%%%%%%%%%%%%%%%%%%%%%%%%%%%%%%%%%%%%%%%%%%%%%
\begin{document}

\begin{titlepage}
    \vspace*{-2.5cm}
    {\centering
    \includegraphics[width=1.00\textwidth]{img/logofiuba.png}\\[2.25 cm]}
    \centering
    \textsc{\Large CB100}\\[0.2 cm]
    \textsc{\large Algoritmos y Estructuras de Datos}\\[4 cm]
    \textcolor{cyan}{{\fontsize{40}{60}\selectfont \bfseries \thetitle}}\\[0.5cm]
    {\Large \bfseries Trabajo Práctico N$^\circ$\tpnumber}\\[5cm]


    \vfill
    % \vspace{10em}
    \noindent\makebox[\linewidth]{\rule{\textwidth}{0.4pt}}\\[0.5cm]
    \begin{minipage}{.46\textwidth}
    \textbf{Autor}\\
    \theauthor
    \end{minipage}%
    \begin{minipage}{.34\textwidth}
    \textbf{Legajo}\\
    \padron
    \end{minipage}%
    \begin{minipage}{.2\textwidth}
     \begin{flushright}
        \textbf{Fecha}\\
        \thedate
    \end{flushright}
    \end{minipage}
\end{titlepage}

%%%%%%%%%%%%%%%%%%%%%%%%%%%%%%%%%%%%%%%%%%%%%%%%%%%%%%%%%%%%%%%%%%%%%%%%%%%%%%%%

{
    \hypersetup{linkcolor=black} % colorlinks=true option is used
    \tableofcontents
    \pagebreak
}

%%%%%%%%%%%%%%%%%%%%%%%%%%%%%%%%%%%%%%%%%%%%%%%%%%%%%%%%%%%%%%%%%%%%%%%%%%%%%%%%

\section{Introducción}

En este trabajo práctico se desarrolla una aplicación de consola que busca
números primos hasta un valor máximo determinado por el usuario, cuando
finaliza, los imprime en un archivo de texto plano en el directorio de ejecución
de la aplicación.

El lenguaje de programacion utilizado para el desarrollo de la aplicación es C++
y el algoritmo utilizado para hallar los números primos es la Criba de
Eratóstenes.

\subsection{Números Primos}

Los números primos por definición son aquellos números positivos que solo son
divisibles por \code{1} y por si mismos, divisibles en términos de que el resto
de la división es nulo. El \code{0} y el \code{1} son dos casos particulares en
los cuales ambos son considerados no primos, el \code{0} por razones obvias, no
está definida la división por \code{0}, y el \code{1} por que solo es divisible
por si mismo, lo cual hace que solo sea divisible por un número.

El número \code{7}, por ejemplo, es un número primo, ya que solo es divisible
por \code{1} y por \code{7}. El número \code{8} no es un número primo, ya que
ademas de ser divisible por si mismo y por \code{1}, es divisible por \code{2} y
por \code{4}.

% \subsection{Algoritmo Primitivo para Descartar Números Primos}
% 
% El primer algoritmo que surge casi instantáneamente es un algoritmo en el que se
% itera sobre una lista de números y por cada numero se descartan todos sus
% múltiplos. Este algoritmo cumple la función de descartar números no primos, pero
% no es para nada eficiente.

\subsection{Criba de Eratóstenes}

La criba de Eratóstenes es un algoritmo sofisticado y eficiente para descartar
numeros no primos de una lista de numeros. En la Criba de Eratóstenes se itera
número a número sobre una lista ordenada y finita de números naturales,
comenzando por el primero, si se trata de un número primo, se buscan todos los
múltiplos en la lista de números y se descartan, luego se avanza al siguiente
número que no haya sido descartado, el cual siempre resulta un número primo, y
se descartan sus múltiplos, así sucesivamente hasta el final de la lista. 

La criba de Eratóstenes se puede optimizar basándose en la propiedad de que los
números no primos (numeros compuestos), se pueden expresar como el producto de
números primos, por ejemplo \code{4 = 2*2} o \code{21 = 3*7}, por lo tanto, si
un número no es primo, debe tener al menos un factor primo que sea menor o igual
que su raíz cuadrada, ya que de lo contrario, el producto de los factores
resultaria en un numero mayor.

De la propiedad anterior se deduce que se puede detener la iteración en la lista
de números cuando se llega a la raiz cuadrado del numero máximo en la lista, ya
que los numeros no primos por encima de éste tendran algun factor menor que la
raiz cuadrada del numero máximo, de este modo se reduce drasticamente el numero
de iteraciones necesarias para descartar los numeros no primos de la lista.

\section{Desarrollo}

Para la implementacion de la Criba de Eratóstenes, como bien se mencionó en la
introducción, se utilizó puramente el lenguaje C++, en particular en su versión
estándar C98. Para organizar el proceso de compilación se utilizó la herramienta
\code{make}.

% Si bien se mencionaron dos algoritmos para hallar números primos, se utiliza
% la Criba de Eratóstenes ya que resulta mucho más eficiente que el algoritmo
% primitivo.

\subsection{Implementación en C++}

Para representar el máximo de números en la lista se define la constante
\code{MAXIMO}.

\begin{fullgrayverb}[\mbox{}]
// Se define la constante `MAXIMO` que representa la cantidad
// de números en la lista
const unsigned int MAXIMO = 100000000;
\end{fullgrayverb}

La Criba de Eratóstenes opera sobre una lista de números naturales ordenada.
Para representarla se utiliza la librería estándar \code{vector} de la cual
se define un vector de tipo \code{bool} inicialmente con valores \code{true}
indicando que son todos números primos, luego al aplicar el algoritmo éste
cambiará el contenido del vector a \code{false} en las posiciones en que se
ubiquen números no primos.

\begin{fullgrayverb}[\mbox{}]
// Se define un vector de tipo `bool` de largo `MAXIMO`
// con valores iniciales `true`
std::vector<bool> numeros(MAXIMO, true);
\end{fullgrayverb}

Comenzando por los casos particulares, \code{0} y \code{1} no son números
primos, por lo que se puede cambiar su valor a \code{false} directamente.

\begin{fullgrayverb}[\mbox{}]
// 0 y 1 no son primos
numeros[0] = numeros[1] = false;
\end{fullgrayverb}

Para determinar los números no primos y descartarlos de la lista de números, se
itera sobre el vector descartando todos los múltiplos de los números primos,
hasta llegar a la raíz cuadrada del numero máximo de la lista. Para determinar
la raíz cuadrada del numero máximo se utiliza la función
\code{std::sqrt} de la librería estándar \code{cmath} 

\begin{fullgrayverb}[\mbox{}]
for (i = 2; i < std::sqrt(MAXIMO); ++i) {
    if(numeros[i] == true) {
        for (j = i; j <= (MAXIMO/i); ++j) {
            numeros[i*j] = false;
        }
    }
}
\end{fullgrayverb}

\pagebreak
Para la impresión en el archivo se define una constante para representar el
nombre del archivo de salida.

\begin{fullgrayverb}[\mbox{}]
// Se define el nombre del archivo a imprimir los números primos
#define OUTPUT_FILE_PATH "primos.txt"
\end{fullgrayverb}

Para la manipulación del archivo se utiliza la clase \code{ofstream} de la
librería estándar \code{fstream}.

\begin{fullgrayverb}[\mbox{}]
// Se define el objeto fp para representar el archivo de salida
std::ofstream fp;
\end{fullgrayverb}

Utilizando los métodos \code{open} e \code{is\_open} del objeto \code{fp} se abre
el archivo, y se comprueba si hubo algún error, en caso afirmativo, se informa
al usuario haciendo uso de métodos de la librería estándar \code{iostream} y se
termina la ejecución con un código de error \code{-1}.

\begin{fullgrayverb}[\mbox{}]
// Se abre el archivo en modo escritura, de haber un error lo reporta
// al usuario y se termina la ejecución inmediatamente
fp.open(OUTPUT_FILE_PATH);
if (fp.is_open() == false) {
    std::cerr << "ERROR: No se pudo abrir `" OUTPUT_FILE_PATH << "`\n";
    return -1;
}
\end{fullgrayverb}

Para imprimir los números primos hallados en el archivo, se recorre la lista en
búsqueda de aquellos casilleros que contengan \code{true}, ya que resultan
números primos, en caso de encontrarse se imprime la posición al archivo, los
casilleros que contienen \code{false} son ignorados. 

La variable \code{j} se utiliza para contar la cantidad de números primos
encontrados, ya que luego será informada al usuario por la consola.

\begin{fullgrayverb}[\mbox{}]
for (i = 2, j = 0; i < numeros.size(); ++i) {
    if(numeros[i] == true) {
        fp << i << std::endl;
        j++;
    }
}
\end{fullgrayverb}

Luego de finalizada la impresión en el archivo, se cierra utilizando el método
\code{close}.

\begin{fullgrayverb}[\mbox{}]
fp.close();
\end{fullgrayverb}

\pagebreak
Finalmente se imprime un mensaje al usuario indicando que terminó la impresión,
y se informa la cantidad de números primos hallados, esto se hace utilizando
metodos de la librería estándar \code{iostream}.

\begin{fullgrayverb}[\mbox{}]
std::cout << "Se encontraron `" << j << "` números primos\n";
\end{fullgrayverb}

\subsection{Makefile}

Si bien es importante la utilización de buenas practicas en el código fuente de
la aplicación, muchas veces el proceso compilación no se tiene en cuenta,
pero es también muy importante para generar una aplicación robusta y fácil de
diagnosticar en caso de fallas.

Para el desarrollo de la aplicación se utiliza la herramienta \code{GNU make},
la cual lee un archivo de nombre exclusivo \code{Makefile} que posee su propia
sintaxis. El archivo \code{Makefile} utilizado intenta ser reutilizable y de
fácil modificación, es por esto que al comienzo del mismo se utilizan variables
para almacenar las principales opciones que el usuario podría querer modificar,
por ejemplo el nombre del compilador utilizado, el nombre final del archivo,
entre otras.

\begin{fullgrayverb}[\mbox{}]
CC := g++
CFLAGS := -Wall -Wshadow -pedantic -ansi -std=c++98 -O3
SRCS := $(wildcard *.cpp)
OBJS := $(SRCS:.cpp=.o)
TARGET := primos
\end{fullgrayverb}

Para compilar y enlazar las cabeceras y el programa final, se utilizan las
siguientes etiquetas:

\begin{fullgrayverb}[\mbox{}]
$(TARGET): $(OBJS)
    $(CC) $(CFLAGS) -o $@ $^

%.o: %.cpp
    $(CC) $(CFLAGS) -c $< -o $@
\end{fullgrayverb}

La etiqueta que contiene \code{\$(TARGET)} se expande al nombre final de la
aplicación, luego de los dos puntos se indican las dependencias, las cuales se
expanden a los archivos objeto (los que terminan en \code{.o})

Para compilar los archivos objetos se utiliza la segunda etiqueta, la cual
corresponde a una \code{pattern rule}, una extension exclusiva de \code{GNU}, 
esta \code{pattern rule} compila todos los archivos terminados en \code{.cpp} a
archivos objeto terminaods en \code{.o}.

\subsection{Optimizaciones del Compilador}

La mayoría de los compiladores de C++ modernos proporcionan una opción para
activar optimizaciones. Al utilizar estas optimizaciones el compilador
intenta eliminar código redundante y optimizar mediante técnicas avanzadas el
archivo binario final, de este modo se logra reducir el uso de memoria, el
tiempo de ejecución y por consiguiente el consumo de energía, entre otras cosas.

En el archivo \code{Makefile} de la aplicación por defecto se utiliza el máximo
de optimizaciones posibles \code{-O3}, de ésta manera se obtiene un archivo
binario superior en rendimiento que el que se obtendría con las optimizaciones
desactivadas.

% Teniendo un compilador que proporcione optimizaciones del código resulta
% conveniente activarla ya que reduce el tiempo de ejecución hasta en un
% \code{406\%}.

\subsubsection{Tiempo de Ejecución}

En el siguiente bloque se muestra la salida del programa utilizando el máximo de
optimizaciones posible del compilador \code{-O3}, se puede apreciar en la salida
del programa el tiempo de ejecución del mismo.

\begin{fullgrayverb}
Se encontraron `5761455` números primos en `11.36` segundos
\end{fullgrayverb}

Desactivando las optimizaciones del compilador, compilando y ejecutando
nuevamente la aplicación, se obtiene el siguiente mensaje en la consola:

\begin{fullgrayverb}
Se encontraron `5761455` números primos en `46.21` segundos
\end{fullgrayverb}

Como se puede observar en el segundo caso, el tiempo de ejecución se incrementa
en un \code{406\%}.

\pagebreak
\section{Instalación}

Para instalar la aplicación usted debe de disponer de una copia del código
fuente, si no posee una, puede obtenerla ingresando al
\href{https://github.com/mjkloeckner/CB100}{Repositorio de Github}, de allí
podrá descargar una copia, o bien puede clonar el repositorio utilizando
\code{git} con el siguiente comando:

\begin{fullgrayverb}
$ git clone https://github.com/mjkloeckner/CB100.git
\end{fullgrayverb}$

Tenga en cuenta que si clona el Repositorio, o lo descarga del repositorio de
Github, tendrá que navegar al directorio \code{tps/1}

\subsection{Sistemas basados en UNIX}

Compruebe que este en el mismo directorio que el archivo \code{main.cpp}. Luego,
compile el código con el programa \code{make}

\begin{fullgrayverb}
 $ make
\end{fullgrayverb}$

Luego puede ejecutar la aplicación de la siguiente manera:

\begin{fullgrayverb}
 $ ./primos
\end{fullgrayverb}$

\subsection{Windows}

En el caso de utilizar Windows como sistema operativo, usted puede configurar
\href{https://en.wikipedia.org/wiki/Windows_Subsystem_for_Linux}{WSL}, el cual
virtualiza un sistema basado en UNIX, para hacerlo puede seguir la
\href{https://learn.microsoft.com/en-us/windows/wsl/install}{guia oficial de
Microsoft}. Una vez configurado WSL siga los pasos para los sistemas basados en
UNIX

De manera análoga si usted dispone de un compilador de C++ puede compilar la
aplicación directamente desde la consola, sin la necesidad de utilizar la
herramienta \code{make}, para eso ejecute el siguiente comando:

\begin{fullgrayverb}
 $ g++ -Wall -Wshadow -ansi -std=c++98 -O3 main.cpp -o primos
\end{fullgrayverb}$

Tenga en cuenta la sintaxis de su compilador ya que puede variar, el comando
anterior está previsto para \href{https://www.mingw-w64.org/}{MinGW-w64}

Luego de compilado la aplicación usted la puede ejecutar de la siguiente manera 

\begin{fullgrayverb}
 $ ./primos
\end{fullgrayverb}$

\section{Conclusión}

Luego de finalizar el desarrollo de la aplicación, se puede concluir que la
combinación de C++ con \code{GNU Make}, resulta muy versátil, y en este caso en
un alto rendimiento.

El lenguaje C++ es uno de los lenguajes de mayor rendimiento en la lista de
lenguajes de programación de alto nivel actuales, junto con C, aunque a
diferencia de éste, C++ ofrece una amplia variedad de librerías estándar, lo que
mejora su robustez y permite al usuario abstraerse de ciertas implementaciones
que resultan irrelevantes en ciertos casos, por ejemplo, en el desarrollo de
esta aplicación, el objecto \code{vector} utilizado para representar la lista de
números naturales. De no existir la librería estándar \code{vector} uno tendría
que ensuciarse las manos e implementar un tipo de dato para representar los
números naturales, ya que no alcanza con los tipos de datos primitivos ofrecidos
por el lenguaje.

Ademas de ahorrar tiempo al programador, las librerías estándar ofrecen mayor
seguridad que si se tratase de un tipo de dato creado por el usuario, ya que las
principales están respaldadas por organizaciones y grupos de estándares como
ISO, IEC, entre otros.

\pagebreak
\section{Referencias}

\begin{itemize}
    \item Deitel H., Deitel P. - C how to Program: With an Introduction to C++
        (2015)
    \item
        \href{https://github.com/mjkloeckner/CB100}{Repositorio de Github}
    \item
        \href{https://en.cppreference.com/w/}{C++ Reference}
    \item
        \href{https://cplusplus.com/reference/}{Standard C++ Library reference}
    \item
        \href{https://en.cppreference.com/w/cpp/container/vector_bool}{
            std::vector\textless bool\textgreater}
    \item
        \href{https://cplusplus.com/reference/fstream/ofstream/}{std::ofstream}
    \item
        \href{https://www.mingw-w64.org/}{MinGW-w64}
    \item
        \href{https://www.gnu.org/software/make/manual/html_node/index.html}{GNU
        Make Manual}
\end{itemize}

\end{document}

El siguiente fragmento del archivo \code{Makefile} utilizado en el desarrollo de
la aplicación intenta ser reutilizable y de fácil modificación, se puede ver al
comienzo del mismo que se definen las principales variables que un usuario
podría querer modificar, por ejemplo el nombre del compilador utilizado o las
flags.
